
% =====================================================================
\section{Non-self-verifiability and the irreducible residue}
\label{sec:residue}
% =====================================================================

The previous sections show the separator has positive content and that this
content silences sub-floor regions. We now reach the deepest obstruction: the
separator cannot be \emph{verified} from within. To verify a part is to decide it
against the totality that constitutes it---its negation-sequence---and that
decision is self-referential, with no consistent value. We prove this by a
self-contained diagonal argument (\Cref{thm:diagonal}), deduce that the
irreducible residue is not a locatable set but a self-referential impossibility
(\Cref{thm:residue-not-measure}), and conclude that the floor can only ever be
bounded below, never exhibited (\Cref{cor:floor-only-bounded}).

\subsection{Verification as deciding membership in the negation}

Recall (\Cref{def:negation}) that $A$ is constituted as $\Space\setminus\Neg(A)$:
its identity is fixed by the totality $\Neg(A)=\Space\setminus A$ of what it is
not. We make verification precise as the decision that constitutes this identity.

\begin{definition}[Verifier]
\label{def:verifier}
A \emph{verifier} for the parts of $\Space$ is a map $\Ver$ that, given a part
$X$ and a candidate constituting totality $T$, returns $\Ver(X,T)\in\{0,1\}$,
intended to read $1$ iff $X$ is exactly the part constituted by withholding $T$,
i.e.\ iff $X=\Space\setminus T$. To \emph{verify $A$} is to evaluate
$\Ver(A,\Neg(A))$: to decide whether $A$ is the part its own negation-sequence
constitutes.
\end{definition}

\begin{remark}[Why verification is membership in the negation]
\label{rem:ver-membership}
By \Cref{def:negation}, $A=\Space\setminus\Neg(A)$ exactly when $A$ and $\Neg(A)$
partition $\Space$; deciding this is deciding, of the constituting totality
$\Neg(A)$, whether what it omits is precisely $A$---a question about $A$ posed
entirely in terms of the collection $\Neg(A)$ that is defined as
``everything but $A$.'' Verification of $A$ is thus inseparable from a question
of self-membership: $A$ figures, through its own definition, inside the totality
against which it is checked.
\end{remark}

\subsection{A self-contained diagonal obstruction}

We now prove, without invoking any external impossibility theorem, that no
verifier decides every part correctly. The argument is a diagonal one, in the
shape of the classical results of~\citet{russell1903,cantor1891,tarski1936} and
the general form of~\citet{lawvere1969}, but it is self-contained: it uses only
the definitions above.

\begin{theorem}[Diagonal obstruction to verification]
\label{thm:diagonal}
There is no verifier $\Ver$ such that, for every part $X$ of $\Space$,
\[
\Ver\bigl(X,\Neg(X)\bigr)=1 .
\]
More precisely, suppose a verifier is \emph{self-applicable}: it can take as its
candidate part the very specification of its own verifying behaviour. Then the
part
\[
D\;:=\;\bigl\{\,\text{parts } X : \Ver(X,\Neg(X))=0\,\bigr\}
\]
\textup{(}the parts the verifier declares not constituted by their own
negation\textup{)} cannot be consistently assigned a verification value:
$\Ver(D,\Neg(D))=1$ and $\Ver(D,\Neg(D))=0$ each entail the other.
\end{theorem}

\begin{proof}
Suppose, for contradiction, a self-applicable verifier $\Ver$ exists. Form $D$ as
stated: $D$ collects exactly those parts $X$ for which the verifier returns $0$
on the self-referential query $\Ver(X,\Neg(X))$. Now ask the verifier about $D$
itself.

If $\Ver(D,\Neg(D))=1$, then $D$ is declared constituted by its own negation;
but by the defining property of $D$, membership-defining condition for $D$ is
$\Ver(X,\Neg(X))=0$, so $D$ satisfying its own membership condition would require
$\Ver(D,\Neg(D))=0$. Reading $D$ through its own definition therefore forces
$\Ver(D,\Neg(D))=0$, contradicting the assumed value $1$.

If instead $\Ver(D,\Neg(D))=0$, then $D$ meets its own membership condition, so
$D$ is among the parts collected by $D$, i.e.\ $D$ is constituted as the part its
negation withholds, which is the assertion $\Ver(D,\Neg(D))=1$, contradicting the
assumed value $0$.

Both values are self-refuting, so no consistent value exists for
$\Ver(D,\Neg(D))$. Hence no self-applicable verifier decides every part; in
particular none returns $1$ on every part, for the construction of $D$ is
available whenever the verifier is total. This is the diagonal obstruction.
\end{proof}

\begin{remark}[Self-application is unavoidable for an embedded agent]
\label{rem:self-app}
The hypothesis of self-applicability is not a technical convenience; it is forced
by embedding. An agent that verifies parts of $\Space$ while being itself a part
of $\Space$ can have its own verifying region presented as a candidate part. The
diagonal part $D$ is then a genuine region, and \Cref{thm:diagonal} applies. An
agent that could stand outside $\Space$ would escape the argument---but no
embedded agent can, and the whole subject of the paper is the embedded agent. The
obstruction is therefore a structural feature of verification-from-within, not an
artefact of an artificial total verifier.
\end{remark}

\subsection{The residue is not a set}

The diagonal obstruction lets us characterise the irreducible residue. One might
hope that the un-verified remainder of an individuation is a definite set---small,
locatable, perhaps measurable---which a finer effort could in principle exhibit
and remove. We show this hope is incoherent: a locatable residue would be
verifiable, and the diagonal obstruction forbids verifying it.

\begin{definition}[Irreducible residue]
\label{def:irreducible}
The \emph{irreducible residue} of the individuation of $A$ is the obstruction to
deciding $\Ver(A,\Neg(A))$: the content that an individuation cannot certify
about itself. It is \emph{exhibitable} if there is a realisable region $E$ with
$\mu(E)>0$ such that verifying $A$ reduces to verifying $A$ on $\Space\setminus E$
together with a separate certification of $E$.
\end{definition}

\begin{theorem}[The irreducible residue is not exhibitable]
\label{thm:residue-not-measure}
Under \Cref{ax:brs} with an embedded \textup{(}self-applicable\textup{)} agent,
the irreducible residue of any individuation is not exhibitable: there is no
realisable region $E$ of positive measure whose separate certification completes
the verification of $A$. The residue is a self-referential impossibility, not a
locatable set.
\end{theorem}

\begin{proof}
Suppose the residue were exhibitable by some realisable $E$ with $\mu(E)>0$.
Certifying $E$ separately means deciding $\Ver(E,\Neg(E))$---verifying $E$ as the
part constituted by its own negation---and, by hypothesis, doing so completes the
verification of $A$. But $E$ is itself a part of $\Space$, and an embedded agent
is self-applicable (\Cref{rem:self-app}); the diagonal construction of
\Cref{thm:diagonal} produces a part on which $\Ver$ has no consistent value, so no
verifier certifies every part, $E$ among them. Hence the supposed certification
of $E$ cannot be guaranteed to terminate with a consistent value; the reduction
fails. Therefore no such $E$ exists: the residue cannot be localised to a
positive-measure set and removed by certifying it. It is the standing
impossibility of the self-referential decision, not a region awaiting finer
treatment.
\end{proof}

\begin{corollary}[The floor can only be bounded below]
\label{cor:floor-only-bounded}
The resolution floor $\thick$ admits a strict positive lower bound
\textup{(}$\thick\ge\mu_{\min}>0$, \Cref{thm:thickness}\textup{)} but no exact
exhibition as $\mu(R)$ for a definite residue set $R$. Every attempt to write the
residue as the measure of a located set $R$ supplies an exhibitable $E=R$,
contradicting \Cref{thm:residue-not-measure}. The floor is known only as the
positive quantity below which no individuation can certify itself.
\end{corollary}

\begin{proof}
The lower bound is \Cref{thm:thickness}. If $\thick=\mu(R)$ for a realisable $R$
of positive measure, then $R$ exhibits the residue, contradicting
\Cref{thm:residue-not-measure}. Hence no such exhibition exists, and $\thick$ is
characterised solely as a positive lower bound on un-certifiable content.
\end{proof}

\begin{remark}[The recognition gap]
\label{rem:recognition-gap}
\Cref{thm:residue-not-measure} has a reading in terms of the relation between an
agent and the part it would verify. To certify $A$ completely, the agent must
decide $A$ against the totality $\Neg(A)$ that constitutes $A$; but the agent,
being embedded, is itself within that totality, so a complete certification would
require it to occupy both sides of the constituting partition at once---to be at
once the part verified and a portion of the totality against which it is
verified. The residue is the irreducible gap that keeps the two apart. Closing
the gap would collapse the distinction on which verification depends; the floor
$\thick$ is the measure of how far the gap can never be closed. This is the
precise sense in which the boundary of knowledge is constitutive, not contingent:
the un-knowable remainder is the very condition under which anything is
individuated at all.
\end{remark}
